\documentclass[a4paper,11pt]{article}

\usepackage{amsmath,amssymb, amsthm}
\usepackage{charter}
%\usepackage{calc}

\usepackage{fullpage}

\usepackage{multirow}
%\usepackage[cp1250]{inputenc}
%\usepackage[T1]{fontenc}
%\usepackage{calligra}
%\usepackage[slovak]{babel}
\usepackage{amsfonts}
\usepackage{layout}
\usepackage[dvips]{graphicx}
\usepackage{url}
\usepackage{color}
\usepackage{graphicx}
\usepackage{algorithmic}
\usepackage{algorithm}
\usepackage{verbatim}
\usepackage{epsfig}
\usepackage[colorlinks,
            linkcolor=red,
            anchorcolor=blue,
            citecolor=blue
            ]{hyperref}
\usepackage{xcolor}

\renewcommand{\algorithmicrequire}{\textbf{Input:}}
\renewcommand{\algorithmicensure}{\textbf{Output:}}

\newcommand{\hlight}[2]{\noindent\colorbox{#1}{%
    \parbox{\dimexpr\linewidth-1\fboxsep}% a box with line-breaks that's just wide enough
        {#2%
        }}}

\newcommand{\HRule}{\rule{\linewidth}{0.5mm}}
 \newcommand{\xtilde}{\tilde x}
 \newcommand{\vtilde}{\tilde v}
 \newcommand{\Embb}{\E}
 \newcommand{\xbar}{\bar x}
 \newcommand{\ie}{\textit{i.e.}}
 
\newcommand{\main}[1]{\footnotesize\textbf{#1}} 
\newcommand{\todo}[1]{{\color{red}#1}}
\newcommand\tagthis{\addtocounter{equation}{1}\tag{\theequation}}


\newcommand{\del}[1]{{\color{red}#1}}

\let\la=\langle
\let\ra=\rangle

\newcommand{\st}{\;:\;}
\newcommand{\ve}[2]{\langle #1 ,  #2 \rangle}

\newcommand{\eqdef}{\stackrel{\text{def}}{=}}

\newcommand{\ii}{{}^{(i)}}

\newcommand{\R}{\mathbb{R}}
\newcommand{\Prob}{\mathbf{Prob}}
\newcommand{\E}{\mathbb{E}}
\newcommand{\Q}{\mathbb{Q}}
\newcommand{\Z}{\mathbb{Z}}

\newcommand{\vc}[2]{#1^{(#2)}}
\newcommand{\nc}[2]{{\color{red}\|#1\|_{(#2)}}}
\newcommand{\ncs}[2]{\|#1\|^2_{(#2)}}
\newcommand{\ncc}[2]{{\color{red}\|#1\|^*_{(#2)}}}
\newcommand{\ls}[1]{{\color{red} \mathcal S(#1)}}
\newcommand{\Rw}[2]{\mathcal R_{#1}(#2)}
\newcommand{\Rws}[2]{\mathcal R^2_{#1}(#2)}
\newcommand{\m}[1]{~\mbox{#1}~}

\newcommand{\nbp}[2]{\|#1\|_{(#2)}}   % norm block primal
\newcommand{\nbd}[2]{\|#1\|_{(#2)}^*} % norm block dual

\newcommand{\lf}{\mathcal L}
\newcommand{\U}{U}
\newcommand{\N}{N}
\newcommand{\mc}[1]{\mathcal #1}

\newcommand{\mLi}{{\color{red}m^{(i)}}}
\newcommand{\gLi}{{\color{red}g^{(i)}}}
%\newcommand{\TLi}{{\color{red}T_L^{(i)}}}
\newcommand{\TLi}[1]{{\color{blue}T^{(#1)}}}

\newcommand{\Lip}{L}

\newcommand{\Rc}[1]{{\color{red}  \mathbf{RC}_{(#1)}}}
\newcommand{\NRCDM}{{\color{red}NRCDM}\  }
\newcommand{\nnz}[1]{{\color{red}\|#1\|_0}}
% sets
\DeclareMathOperator{\card}{card}       % cardinality of a set
\DeclareMathOperator{\diam}{diam}       % diameter of a set
\DeclareMathOperator{\MVEE}{MVEE}       % minim volume enclosing ellipsoid of a set
\DeclareMathOperator{\vol}{vol}         % volume of a set

\DeclareMathOperator{\prox}{prox}         

% statistical
\DeclareMathOperator{\Exp}{\mathbf{E}}           % expectation
\DeclareMathOperator{\Cov}{Cov}         % covariance
\DeclareMathOperator{\Var}{Var}         % variance
\DeclareMathOperator{\Corr}{Corr}       % correlation

% functions and operators
\DeclareMathOperator{\signum}{sign}     % signum/sign of a scalar
\DeclareMathOperator{\dom}{dom}         % domain
\DeclareMathOperator{\epi}{epi}         % epigraph
\DeclareMathOperator{\Ker}{null}        % nullspace/kernel
\DeclareMathOperator{\nullspace}{null}  % nullpsace
\DeclareMathOperator{\range}{range}     % range
\DeclareMathOperator{\Image}{Im}        % image
\DeclareMathOperator{\argmin}{argmin}        % argmin

% topology
\DeclareMathOperator{\interior}{int}    % interior
\DeclareMathOperator{\ri}{rint}         % relative interior
\DeclareMathOperator{\rint}{rint}       % relative interior
\DeclareMathOperator{\bdry}{bdry}       % boundary
\DeclareMathOperator{\cl}{cl}           % closure

% vectors, matrices
\DeclareMathOperator{\linspan}{span}
\DeclareMathOperator{\linspace}{linspace}
\DeclareMathOperator{\cone}{cone}

\DeclareMathOperator{\tr}{tr}           % trace
\DeclareMathOperator{\rank}{rank}       % rank
\DeclareMathOperator{\conv}{conv}       % convex hull
\DeclareMathOperator{\Diag}{Diag}       % Diag(v) = diagonal matrix with v_i on the diagonal
\DeclareMathOperator{\diag}{diag}       % diag(D) = the diagonal vector of matrix D

\DeclareMathOperator{\Arg}{Arg}         % Argument

%\renewcommand{\qedsymbol}{\ding{114}}


\newtheorem{assumption}{Assumption}
\newtheorem{lemma}{Lemma}
\newtheorem{algorithms}{Algorithm}
\newtheorem{theorem}{Theorem}
\newtheorem{proposition}{Proposition}
\newtheorem{example}{Example}
\newtheorem{remark}{Remark}

\theoremstyle{plain}

\newtheorem{prop}[theorem]{Proposition}
\newtheorem{cor}[theorem]{Corollary}
\newtheorem{lem}[theorem]{Lemma}
\newtheorem{claim}[theorem]{Claim}
%\newtheorem{remark}[theorem]{Remark}

\theoremstyle{definition}

\newtheorem{exercise}[theorem]{Exercise}

\newtheorem{rem}[theorem]{Remark}
\newtheorem{que}[theorem]{Question}
\newtheorem{definition}[theorem]{Definition}

 
\begin{document}

\begin{titlepage}
\begin{center}

% Upper part of the page. The '~' is needed because \\
% only works if a paragraph has started.

\includegraphics[width=0.15\textwidth]{lehigh.jpg}~\\[1cm]

\textsc{\LARGE Lehigh University}\\[1.5cm]

\textsc{\Large 418 Final Project}\\[0.5cm]

% Title
\HRule \\[0.4cm]
{ \Large \bfseries On the exact separation of mixed integer knapsack cuts\\[0.4cm] }
%{ \small \bfseries Based on \cite{csizmadia2012s} and some related papers \\[0.4cm] }

\HRule \\[1.5cm]

% Author and supervisor
\noindent
\begin{minipage}{0.4\textwidth}
\begin{flushleft} \large
\emph{Author:}\\
Chenxin Ma, Xi He
\end{flushleft}
\end{minipage}%
\begin{minipage}{0.4\textwidth}
\begin{flushright} \large
\emph{Supervisor:} \\
Prof.~Ted Ralphs
\end{flushright}
\end{minipage}

\vfill

% Bottom of the page
{\Large Fall 2014}

\end{center}
\end{titlepage}

\tableofcontents
\newpage

\section{Introduction} 
Consider a positive integer $n$, letting $b\in\Q$, $a\in \Q^n$, $l\in \{\Q\cup \{-\infty\}\}^n$, $u\in \{\Q\cup \{+\infty\}\}^n$ and $I\subset [n] :=
\{1,...,n\}$. The Mixed Integer Knapsack Set is defined as
\begin{equation}\label{K}
 K=\{x\in \R^n : a^Tx\leq b, l\leq x\leq u, x_i\in \Z, \forall i\in I\}.
\end{equation}

Furthermore, if we have $c\in \Q^n$ and assume $l_i$ is finite for each $i\in [n]$, then the Mixed Integer Knapsack Problem (MIKP) can be 
described as:
\begin{equation}\label{MIKP}
 \max \{c^Tx:x\in K\}.
\end{equation}

We assume that for each $k=1,2,\dots,n$ either $a_k\neq 0$ or $c_k\neq 0$, otherwise, we could remove variable $x_k$ without affecting the problem.

In this report, we present a new branch-and-bound algorithm for MIKP. The methodology that we propose is a linear-programming-based algorithm which exploits dominance conditions. We further make use of lexicographic-domination conditions to eliminate problems with symmetry. One interesting aspect of this approach is that it differs from traditional linear-programming based algorithms by allowing feasible solutions to be pruned during the branching phase.

It might be very difficult to solve MIKP, even by using very effective mixed integer programming solvers such as CPLEX \cite{cplex2005high}. However, the proposed algorithm is shown to be very effective in solving instances of MIKP, much more effectice than CPLEX in fact, both in the amount of time taken to solve problems as by the size of the branch and bound tree explored to find the optimal solution.

In the following content, we will state procedures aiming to solve MIKP
\begin{enumerate}
\item An easy way of identifying unbounded solutions.
\item A way of pre-processing instances of MIKP.
\item The issue of quickly solving the LP-relaxation of MIKP.
\item A simple branch-and-bound algorithm for MIKP.
\item A enhancement of the branch-and-bound algorithm by introducing domination-criteria.
\end{enumerate} 

And finally, we will analyze the computational results of our algorithm and compare it with the general mixed-integer-programing solver CPLEX.

\section{Algorithms for solving MIKP}

\subsection{Infeasible, unbounded, and trivial instances of MIKP}

We aim to use a simple procedure to verify that our MIKP instances are either infeasible, unbounded or trivial, where 'trivial' stands for that our instances are very easy to solve.

Actually, it is very easy to detect the infeasibility of a problem.
\begin{lemma}
 If there is any variable $x_i$ with $i \in [n]$ such that $a_i > 0$ and $l_i = −\infty$, or such that $a_i < 0$ and $u_i = \infty$, then the problem is feasible. 
\end{lemma}

We define the concept of 'efficiency' which can tell us how valuable it is relative to the amount of capacity it uses up in the knapsack constraints.

\begin{definition}
Consider $k\in [n]$ and define
\begin{equation}
e_k=\begin{cases}
c_k/a_k &\m{if} a_k \neq 0,\\
+\infty &\m{if} a_k=0 \m{and} c_k>0,\\
-\infty&\m{if} a_k=0 \m{and} c_k<0.\\
\end{cases}
\end{equation}

We say that $e_k$ is the efficiency of variable $x_k$.
\end{definition}

Furthermore, we also define some status that use them to claim the situation of our problem. Actually, we say that
\begin{table}
\begin{center}
\begin{tabular}{l | c}
potentiator & $(a_k\leq 0, c_k>0, u_k=+\infty)$ or $(a_k\geq 0, c_k<0, l_k=-\infty)$\\
accumulator & $(a_k< 0, c_k=0, u_k=+\infty)$ or $(a_k> 0, c_k=0, l_k=-\infty)$\\
incrementor & $(a_k> 0, c_k>0, u_k=+\infty)$ or $(a_k<0, c_k<0, l_k=-\infty)$\\
decrementor & $(a_k> 0, c_k\geq 0, l_k=-\infty)$ or $(a_k<0, c_k\leq 0, u_k=-\infty)$
\end{tabular}
\end{center}
\caption{Status of a MIKP}
\end{table}

These four definition are very useful, actually, we can obtain the following lemma
\begin{lemma}
If MIKB is feasible and admits a potentiator, then MIKP is unbounded.
\end{lemma}

\begin{lemma}
If MIKB is feasible, and admits an incrementor $x_i$ and a decrementor $x_j$ such that $e_i > e_j$, then MIKP is unbounded.
\end{lemma}

\begin{proposition}
MIKP is unbounded if and only one of the following conditions hold,
\begin{itemize}
\item MIKP is feasible and admits a potentiator $x_j$.
\item MIKP is feasible and admits an incrementor $x_i$ and a decrementor $x_j$ such that $e_i > e_j$.
\end{itemize}
\end{proposition}

Note that even if MIKP is bounded, it may still admit an accumulator. We further define the 'triviality' of MIKP.
\begin{definition}
Consider an instance of MIKP which is feasible and not unbounded. If MIKP has an accumulator, we say that MIKP is trivial.
\end{definition}

Actually, a trivial MIKP can be easily solved by considering the coefficients of the problem. 
\begin{proposition}
Assume that MIKP is feasible and not unbounded. In addition, let $j$ correspond to an accumulator of MIKP. For each $k\in [n]$ such that $k\neq j$ define:
\begin{itemize}
\item $U_k=\begin{cases}
\lfloor u_k \rfloor &\m{if} k\in I,\\
u_k &\m{otherwise.}
\end{cases}$
\item $L_k=\begin{cases}
\lceil l_k \rceil &\m{if} k\in I,\\
l_k &\m{otherwise.}
\end{cases}$
\item $x_k=\begin{cases}
U_k &\m{if}(c_k>0)\m{or}(c_k=0\m{and}u_k<\infty),\\
L_k &\m{if}(c_k<0)\m{or}(c_k=0\m{and}l_k>-\infty),\\
0 &\m{if} c_k=0 \m{and} x_k \m{is free.}
\end{cases}$
\end{itemize}

In addition, with respect to $k=j$, we claim that
\begin{equation}
x_j=\begin{cases}
\max\{\lceil -\frac{\sum_{k\neq j}a_kx_k-b}{a_j}\rceil, l_k\} & \m{if} a_j<0,\\
\min\{\lfloor -\frac{\sum_{k\neq j}a_kx_k-b}{a_j}\rfloor, u_k\} & \m{if} a_j>0.\\
\end{cases}
\end{equation}

Then, we derive that $x$ is well-defined and corresponds to an optimal solution of MIKP.
\end{proposition}

Actually, we can build an algorithm to detect infeasibility, unbounded and find trivial solutions.
\begin{algorithm}
\caption{Detecting infeasibility unbounded and finding trivial solutions}
\label{a1}
\begin{algorithmic}[1] 
\REQUIRE $c,a,b,e,l,u$
\ENSURE $status$
%\STATE $k_0 \leftarrow −1$
%\STATE $k_1 \leftarrow -1; k_2 \leftarrow -1$
\STATE $e^+ \leftarrow -\infty; e^- \leftarrow +\infty$
%\STATE $status \leftarrow$ problem is not unbounded and is not trivial 
\FOR {$i=1$ \TO $n$}
        \IF {$a_i >0 $ \AND $l_i=-\infty$ \OR $a_i <0 $ \AND $u_i=\infty$}
			  \STATE $status \leftarrow$ infeaible
		\ENDIF
		
		\IF {$x_i$ is a potentiator} 
%			  \STATE $k_0 \leftarrow i$
			  \STATE $status \leftarrow$ potentiator
			  \RETURN
		\ELSIF {$x_i$ is an incrementor and $e_i > e^+$}
			  \STATE $e^+ \leftarrow e_i$
%			  \STATE $k_1 \leftarrow i$
		\ELSIF {$x_i$ is an decrementor and $e_i < e^-$}
			  \STATE $e^- \leftarrow e_i$
%			  \STATE $k_2 \leftarrow i$			
	    \ENDIF
\ENDFOR
\IF {$e^+ > e^-$}
		\STATE $status\leftarrow$ incrementor/decrementor pair
		\RETURN
\ELSIF {$e^- = 0$}
	\STATE $status\leftarrow$ accumulator
	\STATE A solution $x$ is given.
\ENDIF
\RETURN

\end{algorithmic} 
\end{algorithm}

\subsection{Preprocessing an instance of MIKP}

In this section we are concerned with reducing an instance of MIKP to another, equivalent instance of MIKP which is easier to solve. A series of procedures for pre-processing an instance of MIKP are now presented. For a thorough introduction to preprocessing see \cite{savelsbergh1994preprocessing}.

\textbf{Test if MIKP is infeasible, trival or unbounded.}  Using Algorithm [\ref{a1}] to test if MIKP is infeasible, trivial or unbounded.  If MIKP is feasible, not trivial and not unbounded, which means it has no potentiators and no accumulators. In addition if variable $x_i$ is an incrementor, and $x_j$ a decrementor, then $e_i ≤ e_j$.

\textbf{Strength bound.} We give a strength bound for our problem. First, we define
\begin{itemize}
\item $U_k=\begin{cases}
+\infty & \m{if} a_k\leq 0,\\
(b-\sum_{i\neq k, a_i>0}a_il_i-\sum_{i\neq k, a_i<0}a_iu_i)/ a_k &\m{otherwise.}
\end{cases}$
\item $L_k=\begin{cases}
-\infty & \m{if} a_k\geq 0\\
(b-\sum_{i\neq k, a_i>0}a_iu_i-\sum_{i\neq k, a_i<0}a_il_i)/ a_k &\m{otherwise.}
\end{cases}$
\end{itemize}

Then, we redefine the stronger bounds of our problem
\begin{equation}
\begin{cases}
u_k=\min\{u_k, U_k\}, l_k=\max\{l_k, L_k\} & \m{if} k\not \in I,\\
u_k=\min\{\lfloor u_k\rfloor, \lfloor U_k \rfloor \}, l_k=\max\{\lceil l_k \rceil, \lceil L_k\rceil\} & \m{if} k\not \in I.\\
\end{cases}
\end{equation}

\textbf{Fix values of variable.} For a given variable $x_k$, we claim
\begin{equation}
x_k=\begin{cases}
u_k &\m{if} a_k\leq 0 \m{and} c_k \geq 0,\\
l_k &\m{if} a_k\geq 0 \m{and} c_k \leq 0.\\
\end{cases}
\end{equation}

After fixing variables as described above, we can substitute out the values in MIKP and obtain a smaller problem with a new right-hand side. In the smaller problem, each variable $x_k$ satisfies either $(a_k >0, c_k>0)$ or $a_k<0, c_k<0$.

\textbf{Complement variables.} For simplicity, we introduce new variable to make sure the lower-bound is always non-negative. Consider a variable $x_k$, and then we set
\begin{equation}
x_k := \begin{cases}
x_k-l_k &\m{if} -\infty <l_k <0,\\
u_k-x_k &\m{if} l_k = -\infty.
\end{cases}
\end{equation}

\textbf{Sort data.} Sort the variables in order of decreasing efficiency. Break first ties if variables are of integer type or not. Break second ties by value of $a_k$.

\textbf{Aggregate variables.} For any given two variables $x_i$ and $x_j$, $i,j\in I$, if $a_i=a_j$, $c_i =c_j$. We aggregate these two variables into a single variable $x_k$ such that $a_k=a_i, c_k=c_i, l_k=l_i+l_j, u_k=u_i+u_j$ and $k\in I$. This procedure will be very helpful later in spending up the branch and bound algorithm.

After the several steps, our finally propose is to reformulate the original MIKP to the following PP-MIKP
\begin{align}\label{proPP}
\max \quad & \sum_{k\in P\cup N} c_kx_k\\
s.t.\quad & \sum_{k\in P\cup N} a_kx_k \leq b\notag\\
& l_k\leq x_k \leq u_k, \forall k\in P\cup N\notag\\
&x_k\in \Z, \forall k\in I.\notag
\end{align}

where $P=\{k: c_k>0\m{and} a_k>0\}$ and $N=\{k: c_k<0\m{and} a_k<0\}$. The PP-MIKP satisfies the following conditions:
\begin{itemize}
\item  PP-MIKP is feasible.
\item  PP-MIKP is not unbounded, and is not trivial.
\item  The variable indices are sorted by efficiency. 
\item  All variables $x_k$ are such that ($a_k >0$ and $c_k >0$) or ($a_k <0$ and $c_k <0$). 
\item For each $k\in P\cup N$, we have $l_k ≥0$.
\item For each $k \in P \cup N$ , all finite bounds are tight; that is, there exists a feasible solution to MIKP which achieves the bound.
\item There are no two identical variables.
\end{itemize}

\subsection{Solving the LP relaxation of PP-MIKP}

In this section, we discuss how to solve the linear programming relaxation of problem (\ref{proPP}).  That is, the problem LP-PP-MIKP.
\begin{align}\label{proLP}
\max \quad & \sum_{k\in P\cup N} c_kx_k\\
s.t.\quad & \sum_{k\in P\cup N} a_kx_k \leq b\notag\\
& l_k\leq x_k \leq u_k, \forall k\in P\cup N\notag
\end{align}

Note that (\ref{proLP}) is nothing more than a linear programming problem. As thus, any Simplex-based linear programming software package would do to solve it. Goycoolea presented a Simplex-liked algorithm, which extends in a simple way Dantzig’s algorithm for solving the linear programming relaxation of bounded, positive coefficient knapsack problems.

\subsubsection{Phase I Algorithm}
\begin{definition}
We say that $x^*\in\mathbb R^{|P|+|N|}$ is tight for (\ref{proPP}) if,
\begin{equation}
\sum_{k\in P\cup N} a_kx^*_k=b.
\end{equation}
\end{definition}
\begin{definition}
$x^*\in\mathbb R^{|P|+|N|}$ is $k-efficient$ for  (\ref{proLP}) if $l_k\leq x^*_k\leq u_k$ and 
\begin{itemize}
\item  $i\in P$ and $i>k$ implies $x_i^*=l_i$.
\item  $i\in P$ and $i<k$ implies $x_i^*=u_i$.
\item  $i\in N$ and $i>k$ implies $x_i^*=u_i$.
\item  $i\in N$ and $i<k$ implies $x_i^*=l_i$.
\end{itemize}
\end{definition}
It can be easy to show that if there exists a tight feasible solution for (\ref{proLP}), then there exists a
tight optimal solution for  (\ref{proLP}). If $x^*$ is tight and k-efficient for (\ref{proPP}), then $x^*$ is an optimal solution of (\ref{proLP}).

The Phase I Algorithm takes as input an instance of (\ref{proLP}), and does one of two things: (a) It proves that the instance is infeasible, or (b) it generates $x$, a k-efficient solution of the instance, having non-negative slack. The algorithm begins by defining $k=\max\{j\in P\cup U: j\in N \m{and} u_j=+\infty \}$, assuming that if the latter set is empty, then $k =-1$. If $k = -1$, it generates the solution $x$, where
\begin{equation}
x_j := \begin{cases}
l_j &\m{if} j\in P,\\
u_j &\m{if} j \in N.
\end{cases}
\end{equation}
If $k > −1$ it generates the solution $x$, for $j\in(P\cup N)\setminus \{k\}$,
\begin{equation}
x_j := \begin{cases}
u_j &\m{if} j\in P \m{and} j<k,\\
l_j &\m{if} j \in P  \m{and} j>k,\\
u_j &\m{if} j\in N \m{and} j>k,\\
l_j &\m{if} j \in N  \m{and} j<k,\\
\end{cases}
\end{equation}
and
$$x_k = \max\Big\{-\frac{1}{a_k}\Big(\sum_{j\neq k} a_jx_j-b, l_k \Big)\Big\}.$$
When $k=-1$, then the algorithm may generate an infeasible solution. In this case it is easy to see that the problem itself is infeasible. Also note that when the algorithm generates a feasible solution, this solution will be efficient. Thus, if the solution is tight it will be optimal. 

\subsubsection{Phase II Algorithm}

The primal phase II algorithm takes as input an efficient solution of LP-PP-MIKP with non-negative slack, and finds from this an optimal solution to the problem. For this, it works by either increasing the values of positive coefficient variables, or decreasing the values of negative coefficient variables in a successive manner until the tightness condition is met, or until all of the variables are at their bounds and the iteration can’t proceed.

\begin{algorithm}
\caption{Primal Phase II Algorithm}
\label{a2}
\begin{algorithmic}[1] 
\REQUIRE $c,a,l,u, P,N, k, x_k, objective, activity.$
\ENSURE $k, x_k, objective, activity, status.$
%\STATE $k_0 \leftarrow −1$
%\STATE $k_1 \leftarrow -1; k_2 \leftarrow -1$
\WHILE{$activity<b$} 
	\IF {$k\in P$}
		\IF{$u_k<+\infty \m{and} a_k(u_k-x_k)<(b-activity)$}
			\STATE $activity\leftarrow activity+a_k(u_k-x_k)$
			\STATE $objective\leftarrow objective+c_k(u_k-x_k)$
		\ELSE 
			\STATE $objective\leftarrow objective+(b-objective)c_k/a_k$
			\STATE $x_k\leftarrow x_k + (b-objective)/a_k$
			\STATE $activity\leftarrow b$
			\STATE $status\leftarrow$ tight optimal
			\RETURN
		\ENDIF
	\ELSE
		\IF{$|a_k|(u_k-x_k)<(b-activity)$}
			\STATE $activity\leftarrow activity+|a_k|(u_k-x_k)$
			\STATE $objective\leftarrow objective+|c_k|(u_k-x_k)$
		\ELSE 
			\STATE $objective\leftarrow objective+(b-objective)|c_k|/|a_k|$
			\STATE $x_k\leftarrow x_k - (b-objective)/|a_k|$
			\STATE $activity\leftarrow b$
			\STATE $status\leftarrow$ tight optimal
			\RETURN
		\ENDIF
	\ENDIF
	\STATE $k\leftarrow k+1$
\ENDWHILE
\STATE $status\leftarrow$ optimal
\RETURN

\end{algorithmic} 
\end{algorithm}
In order to preserve efficiency at each step it iterates by modifying the variables in order of decreasing efficiency. The algorithm always terminates with an optimal solution of the problem. Algorithm (\ref{a2}) shows how a Primal Phase II procedure for LP-PP-MIKP may be implemented.

\subsubsection{Dual Phase II Algorithm}

The dual phase II algorithm takes as input an efficient solution of LP-PP-MIKP with nonpositive slack, and finds from this, an optimal solution to the problem. For this, it works by either decreasing the values of positive coefficient variables, or increasing the values of negative coefficient variables in a successive manner until the tightness condition is met, or until all of the variables are at their bounds and the iteration cannot proceed. In order to preserve efficiency at each step it iterates by modifying the variables in order of increasing efficiency. The algorithm either terminates with an optimal solution of the problem, or a proof that the problem is infeasible. Algorithm (\ref{a3}) shows how a Dual Phase II procedure for LP-PP-MIKP may be implemented.

\begin{algorithm}
\caption{Primal Phase II Algorithm}
\label{a3}
\begin{algorithmic}[1] 
\REQUIRE $c,a,l,u, P,N, k, x_k, objective, activity.$
\ENSURE $k, x_k, objective, activity, status.$
%\STATE $k_0 \leftarrow −1$
%\STATE $k_1 \leftarrow -1; k_2 \leftarrow -1$
\WHILE{$activity>b$} 
	\IF {$k\in P$}
		\IF{$a_k(u_k-x_k)<(activity-b)$}
			\STATE $activity\leftarrow activity-a_k(x_k-l_k)$
			\STATE $objective\leftarrow objective-c_k(x_k-l_k)$
		\ELSE 
			\STATE $objective\leftarrow objective-(objective-b)c_k/a_k$
			\STATE $x_k\leftarrow x_k - (objective-b)/a_k$
			\STATE $activity\leftarrow b$
			\STATE $status\leftarrow$ tight optimal
			\RETURN
		\ENDIF
	\ELSE
		\IF{$u_k<+\infty \m{and} |a_k|(u_k-x_k)<(activity-b)$}
			\STATE $activity\leftarrow activity-|a_k|(x_k-l_k)$
			\STATE $objective\leftarrow objective-|c_k|(x_k-l_k)$
		\ELSE 
			\STATE $objective\leftarrow objective-(objective-b)|c_k|/|a_k|$
			\STATE $x_k\leftarrow x_k + (objective-b)/|a_k|$
			\STATE $activity\leftarrow b$
			\STATE $status\leftarrow$ tight optimal
			\RETURN
		\ENDIF		
	\ENDIF
	\STATE $k\leftarrow k-1$
\ENDWHILE
\STATE $status\leftarrow$ optimal
\RETURN

\end{algorithmic} 
\end{algorithm}


\subsubsection{A branch and bound algorithm for MIKP}

At each step of the branch-and-bound algorithm, we are presented with an instance of PP- MIKP, and an optimal solution to the corresponding linear relaxation, LP-PP-MIKP. This solution is both tight and efficient, and is fully characterized by the values $k,x_k,activity$ and $objective$. If variable k is of continuous type, we know that the current solution is optimal for PP-MIKP. Likewise, if variable k is of integer type, and in addition, $x_k$ takes on an integer value, then the current solution is also optimal for PP-MIKP.

Now, assume that $k$ is of integer type and $x_k$ currently takes on a fractional value, say $f_k$. The simplest branching rule consists in defining two branches: In the first (which we call the down direction), we change the upper bound of variable $x_k$ so that $x_k ≤ \left \lceil{f_k}\right \rceil $. In the other, (which we call the up direction), we change the lower bound of variable $x_k$ so that $x_k ≥\left \lfloor{f_k}\right \rfloor$ .

After branching we would like to avoid solving the new linear programming relaxation of PP-MIKP to optimality from scratch; rather, we would prefer to hot-start the solve from the current solution that we have. This is very easy to do. In the down-direction we round down $x_k$, and in the up-direction we round up $x_k$. Then, activity and objective are updated. If activity increased, we now have an efficient solution with negative slack, so we use the dual phase II algorithm to re-solve. Otherwise, we have an efficient solution with positive slack, so we use the primal phase II algorithm to re-solve.





\iffalse
In this project, we concern two algorithms developed from \cite{fukasawa2011exact}. The first one is to derive strongly valid inequalities for mixed integer knapsack set \eqref{K}, \ie, for a given 
point $x^\star\in \Q^n$, we focus on
\begin{itemize}
 \item identifying an inequality $(\pi,\pi_0)$ which is valid for $\conv(K)$ (equivalent valid for $K$) and violated by $x^\star$.
 \item proving no such inequality exists
\end{itemize}

This is of great practical importance since by obtaining valid inequalities for \eqref{K}, we can also obtain valid inequalities for general Mixed Integer Programming 
(MIP) problems. Suppose $D\in \Q^{m\times n}$ and $d\in \Q^n$, the mixed integer set can be described as
\begin{equation}\label{P}
 P=\{x\in \R^n: Dx\leq d, l\leq x\leq u, x_i\in \Z, \forall i \in I\}.
\end{equation}

Then if $(a,b)$ in \eqref{K} can be obtained as a nonnegative linear combination of row from $(D,d)$, then $P\in K$ and we can find that any inequality
 which is valid for  $K$ will also be valid for $P$. Actually, the valid inequality is called as $knapsack~inequality$ or $knapsack~cut$ of $P$.
 
The basic idea is to transform the problem of getting a valid inequality to a linear separation problem:
\begin{align}
LP_1: \min &\sum_{i=1}^{n} u_i+v_i \notag\\
s.t. \;  & \pi x^k-\pi_0\leq 0,  \quad \forall k\in\{1,2,..,q\} \notag\\
& \pi x^r\leq 0, \quad~ \qquad \forall  k\in\{1,2,..,t\} \notag\\
& \pi x^*-\pi_0 = 1 \notag\\
& \pi+u-v = 0 \notag\\
&u\geq 0, v\geq 0\notag,
\end{align}
where $\{x^1, x^2,...,x^q\}$ and $\{r^1,r^2,..., r^l\}$ represent the extreme points and rays of $conv(K)$. The following steps can be applied to speed up the procedure
\begin{itemize}
 \item Eliminating variables from $LP_1$
 \item Fixing variables and lifting back again
\end{itemize}

The second one is an efficient branch and bound algorithm with both domination branching and reduced-cost bound improvement (KBB) for the Mixed 
Knapsack Integer Problem \eqref{MIKP}. Most modern algorithms for solving Knapsack Problem are based either on branch-and-bound or on dynamic 
programming. However, the most efficient codes seldom make explicit use of Linear Programming. Here we will develop an LP-based branch and bound 
approach, and three strategies are used to speedup solving the problem:
\begin{itemize}
 \item Detecting unbounded solutions and Preprocessing,
 \item Branch and bound (depth-first-search strategy), by using variable reduced-cost information to improve variable bounds at each node of the tree,
 \item Domination (To use the lexicographic and cost dominance to significantly reducing the search space). 
\end{itemize}

The algorithm exploits dominance conditions, which differs from traditional linear programming based algorithms by allowing feasible 
solutions to be pruned during the branching phase. The idea is that feasible solutions will only be pruned if either (a) they are not 
optimal (cost-domination-criteria), or (b) if they are optimal, but somewhere else in the tree it is known that there is another optimal 
solution. In the computational study, we will try to see how  domination plays an role in this algorithm. In fact, according to the literature, the 
proposed algorithm performs quite well in practice, outperforming the general-use mixed integer programming solver CPLEX. 

%In this project, we will focus on cutting plane for large Mixed Integer Programs (MIPs) in the context of Branch and Cut algorithms, 
%which is studied in \cite{fukasawa2011exact}. More specifically, we will study an efficient algorithm (KBB) for the Mixed Integer Knapsack Program 
%(MIKP). 

\section{Research plan}
We will implement the algorithm in \cite{fukasawa2011exact}, which contains:
\begin{itemize}
\item a simple eight-step pre-processing procedure in order to reduce an instance of MIKP to another, equivalent
instance of MIKP which is easier to solve.
\item a depth-first-search branch and bound algorithm which always branches on the unique fractional variable. We use a simple 
linear programming algorithm, a variation of Dantzig’s algorithm \cite{dantzig1957discrete}, which runs in linear time by taking advantage 
of the fact that variables are sorted by decreasing efficiency.
\item using dominance to improve the branch and bound search. In fact, lexicographic and cost dominance allow us to disregard feasible 
solutions that are not the unique lexicographically smallest optimum solution, hence significantly reducing the search space.
\end{itemize}
After implementing it, we will conduct extensive computational experiments. The goal is to reproduce the improvement that the new algorithm 
brings in \cite{fukasawa2011exact}. Moreover, we will try to combine different branching strategies (see \cite{linderoth2005noncommercial}) and 
bound improvement to compare the different versions of the algorithm. If possible, we will take a glimpse to use this knapsack cut generation 
procedure to solve some classes of mixed integer programming problems.
\fi

\bibliographystyle{plain} 
\bibliography{literature}
 
\end{document}
